\section{The Castaway}

The fate of the explorer of consciousness is like that of a Robinson Crusoe who, castaway on an unknown island, must recreate his reality from his own resources.

The following brief, but rich, paragraph from \emph{The Individual and the Becoming of the World} details the requirements to proceed from Stage Two to Stage Three. If we take literally the Traditional teachings that our ordinary life --- if you can call our life in the Cave ``ordinary'' --- is illusion, then all that we naively accepted as obviously true needs to be questioned and discarded. This path is only for those who ``have the courage for an attack on their convictions'' (Nietzsche).

\begin{quotex}
Before going further, we have to point out the necessity that this critical moment of the ideal history of the individual be brought down to and experienced in its depths. Not until he has doubted and denied everything, not until he has cut himself off from everyone, not he until has suffered the unreality of every reality, the uncertainty of every fact, the darkness of every light; not until he has destroyed every support and every haven and has realized the point of the ``great alone'' --- not until the individual can call himself completely that, not until he is an autonomous and self-conscious being. It is this negative act, this absolute tearing of oneself away from whatever used to provide solidity and consistency --- that now makes him be. Besides, according the powerful saying of Stirner, the ``I'' is not everything, but that which destroys everything; through this absolute negativity that tragic principle dawns in man which --- as was clearly seen in Buddhism --- makes him superior to the whole of nature and to the very kingdom of the ``gods''.
\end{quotex}

In the next segment, Evola points out that every experience is somehow one of ``my'' experiences. So that which is constant and unchanging among the changing world phenomena, is the sense of ``I'', whether implicit or explicit. Evola likens it to the Vedic concept of
\emph{ahamkara}, defined as
\begin{quotex}
I-maker. Personal ego. The mental faculty of individuation; sense of duality and separateness from others. Sense of I-ness, ``me'' and ``mine.'' Ahamkara is characterized by the sense of I-ness, sense of mine-ness, identifying with the body, planning for one's own happiness, brooding over sorrow, and possessiveness.
\end{quotex}


So this I presupposes some other experience, that is, the experience the I is conscious of. He relates it to the Vedic teachings on the Observer
... that is never observed. The text follows:


\begin{quotex}
We can specify the place of such an ``I'' as follows. Every experience is inseparably accompanied by the characteristic, implicit or explicit, of being one of my experiences. The self-reference, the \emph{ahamkara} of Vedic metaphysics, is the elementary condition, without which no reality is conceivable, since the only reality that I can concretely speak about is that which, in one way or another, resolves in one of my experiences. Now it is possible to detach this principle of self-reference from the particular contents of experiences in order to fold it back in a certain way onto itself. Thus we have: ``I'' equals ``I'', that is a bare experience, a possession, something simple and ineffable. This bare experience presupposes, in fact and by right, some other experience --- one can say that it is like the cloth from which all the particular experiences are then cut out: here we have that ``seeing that is never seen'', that ``knowing that is never known'', that point of pure centrality that the Upanishads talk about, and in respect to which every particular experience, phenomenon, or thought is a ``\emph{posterius}``, something that comes afterwards and remains at the periphery.
\end{quotex}


Evola then rejects all new age and philosophical teachings that postulate a ``higher self'' that must be reached. There is only the ``I'' and its experiences or representations. Evola then poses this challenge. First he points out that there can only be one such ``I'', and so there cannot be any other ``I''. Therefore, while we are accustomed to regard others as independent centres of awareness, for Evola, the ``others'' are again representations, that is, objects, and not other subjects. This is a crucial challenge made by Evola, and it must be understood and seen, in order to move forward. Evola links the failure to grasp this to sentimentality or a lack of critical reflexion; complete detachment is the only way. Finally, Evola concludes:

\begin{quotex}
This is a point which it is particularly necessary to draw attention to: whoever, either because of moral and sentimental concerns --- that is, reconnected to the preceding stage of natural evidence --- or because of a lack of critical reflexion, does not succeed in extending doubt to the very reality of other subjects and so conceiving them as nothing other than my representations, has not truly brought that detachment, which we wrote about previously, to its limit, and so has not yet perfectly realized the pure essence of the individual. He is not yet ready for the passage to the third stage since those who have first not known how to doubt \emph{everything}, can have absolute certainty of \emph{nothing}.
\end{quotex}


Here again is the reminder to let go of all theories and opinions, in order to base certainty only on what the I can verify for itself. This is the Principle of Induction.

Since the ``I'' cannot be known as an object, Evola next describes what the ``I'' could be. As expected, he does not see the ``I'' as simply the passive observer of phenomena, but relates it to the will-to-power: the I is the centre of will. Therefore, unlike objects --- which are what they are --- the ``I'' must create itself through its activities (``self-positing''). The ``I'' therefore has no static being. Though not explicitly stated here, Evola has in mind the Tao of the Tao Te Ching. Hence, the activity of the ``I'' is the action -- non action (\emph{wei wu wei}) of Taoism.

Evola likens the transition to stage three to the fate of a Robinson Crusoe, who must now create a new world in order to survive. Once all the theories, opinions, beliefs, and creeds have been discarded in stage two, the individual is now in a similar position to the castaway and is also forced to create a new reality:

\begin{quotex}
Therefore, moving onto the third stage, let us immediately say that in it there is an overcoming of the negative point of view connected to the arising of individuality. As one, whom an unfavourable event had cast onto a desert island, followed, after the initial shock, immediately by the will to live, must look for and create the means for a new existence, so the individual, who feels himself by this point alone with himself in the entire extent of the world, can be brought to drag out from his own interior a principle that can secure a new reality beyond the order of appearance and mere representation, in which every thing up to now had to be submerged. This principle is: \textbf{THE POWER OF CONTROL}. The ``I'', in fact, is not a thing, a ``given'', a ``fact'', but, essentially, a deep centre of will and power. As ``I'', Fichte says, it is, only insofar as it posits itself --- and only a pure self-positing is, to tell the truth, its ``being''.
\end{quotex}


Evola now concludes the discussion of the stages before moving on to the relationship between the individual and the becoming of the world.

He points out the the awakening of the I --- or ``still point'' --- in the second stage is now understood in the third stage. This happens not through learning some theory or studying some philosophy, but rather through a process of self-analysis, that is, of one's consciousness. So, what started out (in stage one) as an independent reality is now understood in its relationship to the ``I'' --- both as its representations and also as subject to the will of the I.

Hence, the only principle of certainty that Evola is willing to admit is what can be verified in the consciousness and will of the I.

\begin{quotex}
Through a further self-analysis, the nature of that still point, which was realized in the second stage, is revealed to be such. Now this still point can convey is own substance to whatever lacks it,and that obviously when the various arrangements of that reality, that first appeared irrationally as raw contingency, without participation of the will of the I --- almost as in a dream --- get recovered in accordance with its own connection to an unconditioned assertion of the individual. What now remains is to proceed to a definition of this stage, so that the point of the present treatment is resolved, which is the relationship of the individual to the becoming of the world. In the meantime, we can specify the criterion of certainty that is called for at this point. It is expressed by the principle:

``\emph{There is absolute certainty --- and postulatable reality --- only of those things, of being or non-being, of being such or being otherwise than what the I has in itself, in its capacity of control, or the beginningor cause; of other things, only to the extent of that in them which satisfies such a criterion}``.

These things, which depend in fact entirely on the powerof the ``I'', contribute to the intrinsic evidence that is inherent in its unvarnished beginning.
\end{quotex}



\flrightit{Posted on 2011-02-25 by Cologero}

\begin{center}* * *\end{center}

\begin{footnotesize}
\begin{sffamily}
\texttt{Aidan on 2011-02-25 at 22:44 said:}

Evola errs gravely where Guenon does not. In this essay, we see precisely where and how he departs from Tradition in order to embark upon a way of his own. There is a `tradition' of such departures, of course, with which many identify and self-associate. Gornahoor appears ambiguous in this regard, blending authentic with inauthentic, at times wonderfully accurate, at other times misleading. This is surely because the works of Guenon and Evola are not wholly or even largely compatible.

It is certainly true that one must recognise ahamkara and come to terms with it. Having done so, the point is then to master and eliminate it, not affirm it. In Vedic metaphysics, although it is a relatively important centre, it is by no means ultimate. Indeed, if ahamkara is present, the Self is not. It is astonishing to consider how perverse and self-limiting this makes Evola's version of what is at stake. This is not a minor issue but central and contextual; it is of a piece with his inversion of the relationship between knowledge and action.

The fact that there is a single `I' does not mean there are no other subjects: it means that all other subjects are `I'. There is a vast difference between the pathological spirit of the assumption that others are merely ``representations, that is, objects'' and the view that these `others' are not just the kind of thing *I* am (i.e., other subjects) but are literally `I'. This last perspective alone is compatible with true knowledge since it alone pertains to an order of being that is genuinely metaphysical in nature.

If I regard you as me in spirit, I will not consider you or any other in purely instrumental terms and consequently will not become ensnared in attempts at bending you, them or the world to my will. Such manipulations are paralysingly counter-productive and time-wasting at best. This leaves little to be said for fascist projects, grand totalitarian visions, and all the little `tyrant-creators', spiritual and political, who derive heroic and noble airs from their belief in such systems. Once I know that others are `I', I will then seek to know who I am. Knowing that, I will know who these others are and all other things in essence.

Seeing others as `I' does not deprive these others of their own subjectivity but it does entail that their subjectivity has precisely nothing whatsoever to do with me. What this means is that subjectivity is not mine alone; only MY subjectivity is mine alone: through it, I may find my way (or lose it). There is a universality to this which transcends individuality. One can see how the soul may begin to move and the limitlessness to which it may return.

Far from ending up a castaway on a desert island -- which, incidentally, is a perfect description of the interior isolation that must result from the wrong turn Evola takes -- one who is upon the true path becomes the world and the world becomes a mirror in which the image of the Divine is reflected. THIS is the `great alone'. Not only does it reconcile me with my fellows, as Christ is reported to have instructed me to be prior to approaching God, but, with Christ, I am able to say: ``I and the Father are one.'' Horizontal and vertical are reconciled within the centre and this whole is the Heart.

Contrary to what Evola suggests about the intellectual weaknesses of those who do not follow his way, when all are `I', this Heart is not filled with sentiment; it knows no others, being One alone. Atman is Brahman. That art thou. Is this not the meaning of the Christian injunction: ``be perfect, as your Father in Heaven is perfect''?

\hfill

\texttt{Cologero on 2011-02-26 at 00:36 said:}

Thank you for the thoughtful comment, Aiden. However, you have brought up too many issues to deal with adequately. Nevertheless, I shall try, at least, to deal with the external issues.

As for Gornahoor, we are trying to bring to light the rather complex writings of Traditional writers, particularly Guenon. We hope we are not being ``original'' but rather faithful to Tradition. We also recognize that many more readers have come to Tradition by way of Evola rather than Guenon, so we take some pains to make clear what their respective positions are. We don't necessarily advocate everything we write about; our role is more educational. Perhaps you are not aware that we will be focusing primarily on Guenon's metaphysical works beginning in March. However, that will not be for public consumption and will be restricted to a private group. That ought to make clear to you where our true focus lies.

As for Evola, he is closer to Guenon that you seem to think; both Guenon and Coomaraswamy accepted him, despite some differences. The book in question was written by a young Evola who, having studied German Idealism in some depth, then discovered Guenon. Perhaps the mix of those two strands of thought is not felicitous; but to be generous, perhaps we can give him some credit for attempting to express Traditional ideas in terms that a 20th century European could understand. (By the way, the lectures we are commenting on were delivered to the Theosophical Society in Rome.)

What are the options, Aidan? Both Guenon and Evola conceded the Traditional nature of historical Catholicism, but despaired of the contemporary Church. Guenon opted for Islam, but Islam is an acquired taste. Evola tried to reclaim a Western tradition, even if not totally successfully; his refusal to follow a viable Traditional path in practice may account for some of his anomalies. By the way, Evola did write an interesting essay on what an authentic esoteric Catholicism would look like. We are working on translating that piece into English, but again, just for private circulation.

As for the facile and misleading accusation regarding Fascism, we have written about that topic more than once. I suggest you look it up and ponder it.

As for the actual content on Evola's argument: the task is more subtle than you make it sound. There is the epistemological question: how do I know other I's even exist? Presumably you know it by inference, but that is not what is meant by gnosis. Is the inability to know other I's sufficient warrant to answer the ontological question about the existence of other I's in the negative? You claim there are multiple I's; does that imply then that there are multiple Brahmans?

Also, please don't jump to the conclusion that Evola's position is egoism. There is no question here of using others instrumentally; as a matter of fact, Evola in his other writings insisted on free and organic relationships among men. If there is one ``I'' there is no ``other''; a fortiori, there is not even an ``ego'' with grand projects, etc. To the contrary, the other represents privation to me, where my will doesn't extend to. Privation is a concept derived from Guenon and is a topic we shall deal with soon. It is similar to the Kabbalistic notion of tsimtsum: viz, God voluntarily withdraws to create a clearing for others to manifest.

In this context, you may want to read Guenon's essay ``Is the Spirit in the Body or the Body in the Spirit?''. We can't delve into it in a comment, but perhaps as a post some day. But here is a taste:

\begin{quotex}

It is indeed the spirit (Atma) that is truly the universal center containing all things; but when reflected in human manifestation it there by appesrs as localized at the center of the individuality, and even, at the center of its corporeal modality... This reflection is assuredly only an appearance, as is individual manifestation itself; this appearance is reality for it .. It is only when a being has gone beyond these limits that the other point of view becomes real for it; its center is thenceforth in the universal.
\end{quotex}


Reread Evola in the light of Guenon and you may get a different impression.

\hfill

\texttt{Aidan on 2011-02-26 at 20:53 said:}

I appreciate your reply, Cologero, and assure you that my enthusiasm and admiration for the treasures to be found at Gornahoor will remain undiminished, regardless of how you position yourself with regard to the criticisms I have made of Evola.

I do see more than a little value in the work of Evola but confess to having concluded that his main value is as a negative example. He is an effective symbol of a way that it is better not to go. His diagnosis of the ills of modernity is powerful and accurate but we might well expect it to be, given the cue he takes from Guenon with respect to the roots of the degeneracy. It is what he decides to do about it and, indeed, how he frames things from there on, that serves to demonstrate the errors in his understanding.

The more we advance -- the higher the level we reach -- the more serious our errors become. As others have noted, Evola's placement of the regal above the sacerdotal is non-traditional, to say the least, and it is difficult to see how it is not fully anti-traditional in its reinforcement of the `revolt of the kshatriyas' (a Luciferic event, if ever there was one). This undermining of brahmin spiritual authority is not a side issue; it sets the tone for what follows and is a recipe for disaster that has been followed by a variety of individuals and groups. It does not merely undermine the higher order but fully and deliberately usurps it. It is this keynote element in his work that allows so many politicals to seize upon Evola as `inspiration' for their activities. If this usurpation is denied legitimacy, as it should be, then there is a basis for individual and collective tranquility, which is the only foundation there can be for the establishment of true order.

I am unsure why you describe my comment about fascism as an accusation. Who do you believe I have accused? I had already read some of what you have written about fascism and, since you suggest it, have reread certain essays. I accept that Evola would have considered himself above fascism and indeed all political movements, identifying himself instead with permanent principle. He was involved in the movement, however, to one degree or another (reports of the exact degree vary wildly), and that he harboured high hopes for European Fascism cannot be denied. Coomaraswamy apparently explicitly rejected it at the time and Guenon had no involvement with it of any kind. This alone entitles fascist movements of whatever stripe to claim him as one of their own and his writings do give comfort to them. When it comes to Guenon, however, they have no legitimate claim, whatever they may think; nothing in his work endorses their understanding of the world. On the contrary, it reveals the limitations and deformities of that understanding, exactly as it does with every other non-traditional programme of thought and action.

As to the options, the ones that are available depend upon the nature of the person. Once that nature is made visible, the appropriate option will also stand revealed. How could it be otherwise?

Finally, I don't know where you found me claiming that there are multiple I's. Did I not agree, and repeatedly state, that there is one `I'? What I did was refer to multiple subjects, multiple others, each and every one referring to `I'. Indeed, this is why I made the above comment in the first place, to highlight the distinction between the subject and the `I' to which it refers, and thereby demonstrate the element missing in Evola's presentation. There is an entire level missing, namely that of the soul. He is -- and you are -- conflating soul (subject) and Self (I).

Contrary to what Evola is saying, the existence of multiple subjects is self-evident. I am replying to you, am I not? I am assuming you also have family, friends, colleagues, correspondents, etc., each of whom is an other. I have already agreed that these others are not other I's -- to think they were would be a mistake and, yes, a privation to the soul of one who took such a view (not because my will does not extend to them but because my being does not).

The fact that all subjects are I does not make them objects, however -- indeed, it makes them plainly NOT objects. The difference between a subject and an object is of immense significance on a variety of levels but it is too much to go into in further detail here. Have I made it all sound too easy and straightforward? The existence of other subjects is an irreducibly simple matter; by contrast, the knowledge that all are I is beyond subtle. Everyone, apart from a few philosophers, knows the former; who knows the latter?

Although concerned at having been overly bold, I had believed my earlier comment to be quite clear. I hope I have now removed any confusions that may have been the result of a lack of clarity on my part.

\hfill

\texttt{HOO on 2011-02-27 at 05:09 said:}

Where or when exactly does Evola undermine brahmin spiritual authority? Where does he put *mere* warriors above those initiated?

\hfill

\texttt{Aidan on 2011-02-27 at 12:41 said:}

HOO, it is not a question of placing ``*mere* warriors above those initiated''. Evola looks to a REGAL initiation so sublime and glorious that it subsumes the sacerdotal function within itself (prior to the delegation of that function to those of a lesser state, i.e., priests). He looks to a noble elite worthy of being initiated into a heroic warrior caste designed to bring about a ``revolution from above''.

If you doubt the words I am personally placing upon all this and think me guilty of distortion, might I offer you this address?

\url{http://thompkins_cariou.tripod.com/id95.html}


It should lead you to a fine site called `Evola As He Is' and an essay by Evola called, appositely, ``Spiritual Authority and Temporal Power''. If you accept the argument he lays out -- which expresses the general spirit of his work very well -- perhaps you will not be inclined to see it as `undermining' or `usurping' but as a simple assertion of something true and correct. It is worth considering it from the other side, however; that is, reading Evola in the light of Guenon, as Cologero advised me to do; in other words, begin with the reverse of the proposed inversion. If Guenon is correct, what are the consequences of such a profound inversion of the truth likely to be, especially when it occurs at a level that is high but not the highest?

The divine order exists everywhere eternally in principle, even if it exists nowhere in practice for a particular time. It is certainly frustrating to live in a society in which the divine order does not exist in practice. On the other hand, what else is to be expected under the conditions of Kali Yuga? The hidden truth is that the degeneracy we see all around us is in order and must be endured. This does not mean we are to participate in its downwardness, or facilitate the spread of it. Whatever the Age, it is the task of all castes, natures and types of men to conform to the divine order in their own way, in order that it may exist in them by reflection. This is the Way and there is no other way. Every action that does not disturb this reflection is legitimate. Within the domain of such legitimacy -- which is vast but not limitless -- we find the noble virtues of the kshatriyas united within and beneath the regal function; this is the true solar path.

Since the ordering is obvious for those with eyes to see it, what Evola proposes is plainly a type of disorder and a serious one at that. Although it is tempting to do so, it would be unfair to sink beneath logic and seize upon the course and temporal consequences of European Fascism as 'empirical' evidence of the error, especially since Evola took both regimes to task for not being right-wing enough (hence his ability to title one of his books `Fascism Seen From The Right'). He could distinguish himself from the fascist movement -- I do not doubt that his own vision was far nobler and higher than anything to which the majority of its adherents would have aspired. Most fascists would have a hard time distinguishing themselves from him, however. For them, he is likely to remain an intellectual high water mark.

\hfill

\texttt{Cologero on 2011-02-27 at 12:43 said:}

Regarding royalty. Guenon writes in ``Lord of the World'':

\begin{quotex}

By the Middle Ages, supreme power had already become divided between the Empire and the Papacy. SUch a division marks an organization that is incomplete at tis head since the common principle, on which the two powers depends, is missing... In ancient Rome, on the other hand, the Emperor was alos Pontifex Maximus. The Muslim theory of the Caliphate also unites the two powers... \end{quotex}


So Guenon and Evola agree on this issue. We wrote about this in the \textit{Caste and Social Order}\footnote{\url{http://www.gornahoor.net/?p=1698}} and also on the \textit{City of the Sun}\footnote{\url{http://www.gornahoor.net/?p=518}}. So it is certainly a traditional notion.

As for Evola's undermining of spiritual authority, please post a specific reference so we can discuss.

So you and Evola agree about the singular ``I''. I don't understand the objection.

To repeat myself again. Evola rebutted the post WWII accusation of fascism by pointing out that his position is what every sane and well-bred man believed before 1789. Guenon's position, although he avoided political involvement, would not be much different. Nevertheless, Guenon maintained relationships with leaders of Action Francaise and allowed some articles to be published in Fascist journals. Unfortunately, the Traditional worldview looks like the caricature of Fascism that everyone seems to believe. Specifically, any nostalgia to ``throne and altar'' is today rejected as ``Fascist'' (although it is not). That is why Guenon is falsely associated with Fascism, although that would be his position.

Again, we must repeat ourselves: different perspectives are appropriate to different castes. Guenon writes as a Brahmin, cold and austere. His goal is complete transcendence; hence his retirement to a life as an Egyptian peasant, where he lived with the idea of the Caliphate, Sharia Law and an esoteric Sufi Order, things he could not find in Europe. Most types are left cold by Guenon.

Evola writes as a Kshatriya, with passion and vigour. His goal is the True Man. Unlike a world-rejecting Guenon, his task is precisely involvement with the world; it is his Nature. As such, he must deal with the world as he finds it, not as he imagines or wishes it to be. It is always a dirty step to apply principles --- which are absolutely true --- to specific situations, which are necessarily ambiguous, dark, and uncertain. Give the guy some slack.

\hfill

\texttt{Will on 2011-02-27 at 12:44 said:}

I don't think it's accurate to say that Evola championed the Ksatriya caste above the Brahmin. He identified himself with the Ksatriya caste, but what he championed was a notion of Regality which does not fit neatly into the Vedic caste system. It seems to be related to the Indo-European tradition of the King-Magus, as elaborated in the works of Georges Dumezil. He wanted a ruler who embodied both temporal and spiritual power -- that was his idea of Ghibellinism. I interpret his concept of solar regality as non-dualist -- it rejects the duality of action and contemplation.

One can see this in his own life. Evola was the least bookish of the major traditionalist writers, opting instead for mountain climbing, military service, and leisurely strolls through war zones. All of that in addition to being incredibly learned. He seems to have attempted to embody the virtues of both action and contemplation in his life. How successful he was is perhaps open to debate, but it seems to me a noble attempt.

The ancient Greek tradition of dual kingship is relevant here. The Spartans kept this tradition up until the time of the Peloponnesian War or so. While Sparta is mostly remembered for its military prowess, it should be noted that they were also thought of as the most religious of all the Greeks. The Spartan ideal of embodying both virility and wisdom is what Evola championed.

I'm not sure how that would fit into the traditional caste system, or if it does at all. But I suspect Evola's fondness for Buddhism is a factor here as well. The Buddha was a ksatriya, but by attaining supreme enlightenment, transcended the caste system altogether. Evola wrote that there were two ways to be outside the caste system: as a pariah, lower than even the servile caste, or as a transcender, higher than even a Brahmin. The Buddha was the latter, and was called a cakravartin and a universal monarch.

So to say that Evola's views are anti-traditional seems to go with Guenon's original position that the whole Buddhist tradition was anti-traditional. But of course he later changed his mind about this, thanks in large part to Ananda Coomaraswamy and, if I remember right, Marco Pallis.

\hfill

\texttt{Aidan on 2011-02-27 at 13:57 said:}

The issue is not whether the two functions -- royal and sacerdotal -- can be united within a single Person. As you point out, they can be and have been. The issue is the relationship of the two functions to the original principle within which they are essentially undivided and, consequently, their relationship to each other in an apparently divided state. Only by the light of this first principle can we see whether the different functions are being discharged correctly in practice.

I don't think I can be any plainer than I have been so am happy to withdraw into silence at this point. There is no sense in repeating ourselves.

Thanks for the observations, Will. I would only add that few if any of the great teachers were still warriors once they had become great teachers, Socrates included. This is because they were now teachers and their function had changed. (If anyone cites the Bhagavad-Gita in support of Evola's world-order -- as Evola himself did -- note that it is Arjuna, the Kshatriya, who is under instruction.) If you transcend the caste system, you cease to be of the caste from which you came. In Vedic terms, those who stand outside caste due to having transcended all conditions are known as jivan-mukti, of whom Sri Ramana Maharshi is a notable Twentieth Century example. These ones are capable of teaching the teachers. Since they are in themselves the fulfilment of that toward which the brahmin is oriented, it is reasonable to look to them for ways to counteract the effects of brahmin degeneracy.

Forgive me if I now withdraw from further comment. I do not wish to draw you away from your purpose at this site and have my own duties to which I must attend. Many thanks for the depth of consideration given to my possibly impertinent remarks.

\hfill

\texttt{HOO on 2011-02-27 at 15:29 said:}

Often seeen and predictable Guénonianist prejudice and veneer.

´The brahmaa caste is habitually thought of in the West as a ``sacerdotal'' caste. This is true only up to a certain point. In the Vedic origins the type of Brahman or ``sacrificer'' bears little resemblance to that of the ``priest'' as our contemporaries think of him: he was, rather, a figure both virile and awful and, as we have said, a kind of visible incarnation in the human world of the superhuman (bhu-deva). Furthermore, we often find in the early texts a point where the distinction between the brahma---the ``sacerdotal'' caste---and the ksatram or rajam---the warrior or regal caste---did not exist; a feature that we see in the earliest stages of all traditional civilizations, including the Greek, Roman, and German. The two types only began to differ in a later period, this being another aspect of the process of regression that we have mentioned.´ Evola, ``The Doctrine of Awakening'' p. 28 \hfill

\texttt{kadambari on 2011-03-03 at 09:05 said:}

HOO,

Great comment, you have some insight. As you can see, most people bandy the terms ``Vedic caste system'' without a clue as to its origins or what it stood for and what it means... .

\end{sffamily}
\end{footnotesize}

